\documentclass[../../main]{subfiles}

\begin{document}

\section{Fundamentos}
\label{sec:matematica:funcoes:fundamentos}

\begin{defn}[Função]
	\label{def:matematica:funcoes:fundamentos:funcao}

	Sejam (A) e (B) conjuntos.

	Uma função de (A) em (B) é uma relação

	[
	f \subseteq A \times B
	]

	que satisfaz as seguintes propriedades:

	\begin{enumerate}

		```
		\item Para todo \(a\in A\), existe um elemento
		      \(b\in B\) tal que

		      \[
			      (a,b)\in f.
		      \]

		\item Para todo \(a\in A\) e para quaisquer
		      \(b,c\in B\),

		      \[
			      (a,b)\in f
			      \land
			      (a,c)\in f
			      \Rightarrow
			      b=c.
		      \]
		      ```

	\end{enumerate}

\end{defn}

\begin{obs}

	A primeira condição é denominada
	\emph{existência da imagem}.

	A segunda condição é denominada
	\emph{unicidade da imagem}.

	Em conjunto, essas propriedades garantem que cada
	elemento do domínio possui exatamente uma imagem.

\end{obs}

\begin{notation}

	Quando (f) é uma função de (A) em (B),
	escreve-se

	[
	f:A\to B.
	]

\end{notation}

\begin{prop}[Caracterização Funcional]
	\label{prop:matematica:funcoes:fundamentos:caracterizacao}

	Uma relação

	[
	f\subseteq A\times B
	]

	é uma função se, e somente se, para cada
	(a\in A) existe exatamente um elemento
	(b\in B) tal que

		[
			(a,b)\in f.
		]

\end{prop}

\begin{proof}

	A condição de existência fornece pelo menos uma
	imagem para cada elemento de (A).

	A condição de unicidade garante que essa imagem é
	única.

	Reciprocamente, a existência de exatamente um
	elemento (b\in B) para cada (a\in A)
	implica imediatamente as duas propriedades da
	definição.

\end{proof}

\begin{defn}[Domínio]
	\label{def:matematica:funcoes:fundamentos:dominio}

	Seja

		[
			f:A\to B.
		]

	O conjunto (A) é denominado domínio de (f).

	Escreve-se

	[
	\operatorname{Dom}(f)=A.
	]

\end{defn}

\begin{defn}[Contradomínio]
	\label{def:matematica:funcoes:fundamentos:contradominio}

	Seja

		[
			f:A\to B.
		]

	O conjunto (B) é denominado contradomínio de
	(f).

\end{defn}

\begin{defn}[Imagem de um Elemento]
	\label{def:matematica:funcoes:fundamentos:imagem-elemento}

	Seja

		[
			f:A\to B.
		]

	Para cada elemento (a\in A), o único elemento
	(b\in B) tal que

		[
			(a,b)\in f
		]

	é denominado imagem de (a) por (f).

	Escrevemos

		[
			f(a)=b.
		]

\end{defn}

\begin{prop}[Unicidade da Imagem]
	\label{prop:matematica:funcoes:fundamentos:unicidade-imagem}

	Seja

		[
			f:A\to B.
		]

	Para cada (a\in A), existe exatamente um
	elemento (b\in B) tal que

		[
			f(a)=b.
		]

\end{prop}

\begin{proof}

	Segue diretamente da definição de função.

\end{proof}

\begin{defn}[Imagem de uma Função]
	\label{def:matematica:funcoes:fundamentos:imagem}

	Seja

		[
			f:A\to B.
		]

	A imagem de (f) é o conjunto

		[
			\operatorname{Im}(f)
				=
				{
					f(a)
					\mid
					a\in A
				}.
		]

\end{defn}

\begin{prop}
	\label{prop:matematica:funcoes:fundamentos:imagem-subconjunto}

	Para toda função

	[
	f:A\to B,
	]

	vale

		[
			\operatorname{Im}(f)
			\subseteq
			B.
		]

\end{prop}

\begin{proof}

	Pela definição de função,
	(f(a)\in B) para todo (a\in A).

	Portanto todo elemento da imagem pertence ao
	contradomínio.

\end{proof}

\begin{defn}[Gráfico de uma Função]
	\label{def:matematica:funcoes:fundamentos:grafico}

	Seja

		[
			f:A\to B.
		]

	O gráfico de (f) é o conjunto

		[
			G_f
				=
				{
					(a,f(a))
					\mid
					a\in A
				}.
		]

\end{defn}

\begin{prop}
	\label{prop:matematica:funcoes:fundamentos:grafico-subconjunto}

	Para toda função

	[
	f:A\to B,
	]

	vale

		[
			G_f
			\subseteq
			A\times B.
		]

\end{prop}

\begin{proof}

	Se

	[
	(a,f(a))\in G_f,
	]

	então (a\in A) e (f(a)\in B).

	Logo

		[
			(a,f(a))
			\in
			A\times B.
		]

\end{proof}

\begin{ex}

	Considere

	[
	f:\mathbb N\to\mathbb N,
	\qquad
	f(n)=n^2.
	]

	Seu gráfico é

	[
	G_f
		=
		{
			(0,0),
			(1,1),
			(2,4),
			(3,9),
			\ldots
		}.
	]

\end{ex}

\begin{prop}[Princípio da Extensionalidade para Funções]
	\label{prop:matematica:funcoes:fundamentos:extensionalidade}

	Sejam

		[
			f:A\to B
			\quad\text{e}\quad
			g:A\to B.
		]

	Então

	[
	f=g
	]

	se, e somente se,

	[
			(\forall a\in A)
			\bigl(
			f(a)=g(a)
			\bigr).
		]

\end{prop}

\begin{proof}

	Se (f=g), então seus gráficos coincidem e,
	consequentemente,

	[
			f(a)=g(a)
		]

	para todo (a\in A).

	Reciprocamente, suponha que

		[
			f(a)=g(a)
		]

	para todo (a\in A).

	Então os pares ordenados

		[
			(a,f(a))
		]

	e

		[
			(a,g(a))
		]

	coincidem para todo (a\in A).

	Portanto os gráficos de (f) e (g) são iguais.

	Logo

		[
			f=g.
		]

\end{proof}

\end{document}
